\documentclass[10pt,a4paper]{amsart}
\usepackage{amsmath,amssymb,amsthm,graphicx}
\usepackage{verbatim}
\newenvironment{reference}{\comment}{\endcomment}
%\newenvironment{reference}{}{}
\newenvironment{slogan}{\comment}{\endcomment}
\newenvironment{history}{\comment}{\endcomment}

% For commutative diagrams we use Xy-pic
\usepackage[all]{xy}

% We use 2cell for 2-commutative diagrams.
\xyoption{2cell}
% Theorem environments.
%
\theoremstyle{plain}
\newtheorem{theorem}[subsection]{Theorem}
\newtheorem{proposition}[subsection]{Proposition}
\newtheorem{lemma}[subsection]{Lemma}

\theoremstyle{definition}
\newtheorem{definition}[subsection]{Definition}
\newtheorem{example}[subsection]{Example}
\newtheorem{exercise}[subsection]{Exercise}
\newtheorem{situation}[subsection]{Situation}

\theoremstyle{remark}
\newtheorem{remark}[subsection]{Remark}
\newtheorem{remarks}[subsection]{Remarks}

\numberwithin{equation}{subsection}

% Macros
%
\def\lim{\mathop{\mathrm{lim}}\nolimits}
\def\colim{\mathop{\mathrm{colim}}\nolimits}
\def\Spec{\mathop{\mathrm{Spec}}}
\def\Hom{\mathop{\mathrm{Hom}}\nolimits}
\def\Ext{\mathop{\mathrm{Ext}}\nolimits}
\def\SheafHom{\mathop{\mathcal{H}\!\mathit{om}}\nolimits}
\def\SheafExt{\mathop{\mathcal{E}\!\mathit{xt}}\nolimits}
\def\Sch{\mathit{Sch}}
\def\Mor{\mathop{\mathrm{Mor}}\nolimits}
\def\Ob{\mathop{\mathrm{Ob}}\nolimits}
\def\Sh{\mathop{\mathit{Sh}}\nolimits}
\def\NL{\mathop{N\!L}\nolimits}
\def\CH{\mathop{\mathrm{CH}}\nolimits}
\def\proetale{{pro\text{-}\acute{e}tale}}
\def\etale{{\acute{e}tale}}
\def\QCoh{\mathit{QCoh}}
\def\Ker{\mathop{\mathrm{Ker}}}
\def\Im{\mathop{\mathrm{Im}}}
\def\Coker{\mathop{\mathrm{Coker}}}
\def\Coim{\mathop{\mathrm{Coim}}}

% Boxtimes
%
\DeclareMathSymbol{\boxtimes}{\mathbin}{AMSa}{"02}

%
% Macros for moduli stacks/spaces
%
\def\QCohstack{\mathcal{QC}\!\mathit{oh}}
\def\Cohstack{\mathcal{C}\!\mathit{oh}}
\def\Spacesstack{\mathcal{S}\!\mathit{paces}}
\def\Quotfunctor{\mathrm{Quot}}
\def\Hilbfunctor{\mathrm{Hilb}}
\def\Curvesstack{\mathcal{C}\!\mathit{urves}}
\def\Polarizedstack{\mathcal{P}\!\mathit{olarized}}
\def\Complexesstack{\mathcal{C}\!\mathit{omplexes}}
% \Pic is the operator that assigns to X its picard group, usage \Pic(X)
% \Picardstack_{X/B} denotes the Picard stack of X over B
% \Picardfunctor_{X/B} denotes the Picard functor of X over B
\def\Pic{\mathop{\mathrm{Pic}}\nolimits}
\def\Picardstack{\mathcal{P}\!\mathit{ic}}
\def\Picardfunctor{\mathrm{Pic}}
\def\Deformationcategory{\mathcal{D}\!\mathit{ef}}
\usepackage{cleveref}
\crefname{originaldefinition}{original definition}{original definitions}
\Crefname{originaldefinition}{Original definition}{Original definitions}
\makeatletter
\newlabel{definition-perfection}{{046X}{source}{}{sourcepointer-1}{}}
\newlabel{definition-perfection@cref}{{[originaldefinition][0][]046X}{[1][0][]source}{}{}{}}
\makeatother
\crefname{originallemma}{original lemma}{original lemmas}
\Crefname{originallemma}{Original lemma}{Original lemmas}
\makeatletter
\newlabel{lemma-finite-type-field-Jacobson}{{00G1}{source}{}{sourcepointer-2}{}}
\newlabel{lemma-finite-type-field-Jacobson@cref}{{[originallemma][0][]00G1}{[1][0][]source}{}{}{}}
\makeatother
\crefname{originallemma}{original lemma}{original lemmas}
\Crefname{originallemma}{Original lemma}{Original lemmas}
\makeatletter
\newlabel{lemma-flat-fibres-irreducible}{{00I6}{source}{}{sourcepointer-3}{}}
\newlabel{lemma-flat-fibres-irreducible@cref}{{[originallemma][0][]00I6}{[1][0][]source}{}{}{}}
\makeatother
\crefname{originallemma}{original lemma}{original lemmas}
\Crefname{originallemma}{Original lemma}{Original lemmas}
\makeatletter
\newlabel{lemma-limit-argument}{{00I3}{source}{}{sourcepointer-4}{}}
\newlabel{lemma-limit-argument@cref}{{[originallemma][0][]00I3}{[1][0][]source}{}{}{}}
\makeatother
\crefname{originallemma}{original lemma}{original lemmas}
\Crefname{originallemma}{Original lemma}{Original lemmas}
\makeatletter
\newlabel{lemma-p-ring-map}{{0BRA}{source}{}{sourcepointer-5}{}}
\newlabel{lemma-p-ring-map@cref}{{[originallemma][0][]0BRA}{[1][0][]source}{}{}{}}
\makeatother
\crefname{originallemma}{original lemma}{original lemmas}
\Crefname{originallemma}{Original lemma}{Original lemmas}
\makeatletter
\newlabel{lemma-perfect-reduced}{{00I4}{source}{}{sourcepointer-6}{}}
\newlabel{lemma-perfect-reduced@cref}{{[originallemma][0][]00I4}{[1][0][]source}{}{}{}}
\makeatother
\crefname{originallemma}{original lemma}{original lemmas}
\Crefname{originallemma}{Original lemma}{Original lemmas}
\makeatletter
\newlabel{theorem-nullstellensatz}{{00FV}{source}{}{sourcepointer-7}{}}
\newlabel{theorem-nullstellensatz@cref}{{[originallemma][0][]00FV}{[1][0][]source}{}{}{}}
\makeatother
\title{Over a separably closed field the tensor product of two algebras with a unique minimal prime has one too}
\author{The Stacks Project Authors}
\date{stacks-project:algebra.tex:local-snapshot-20261012; GNU FDL 1.2 or later}
\begin{document}
\maketitle
\noindent\textit{Editorial note.} All delivered mathematical text is the author's own, copied from the pinned Stacks Project chapter file algebra.tex; no statement or proof step was rewritten, reordered or removed.\par\medskip
\noindent\textit{Editorial note.} Preamble provenance, as required for this source: the frozen snapshot stages the chapter's own preamble file beside it (preamble.tex, 217 lines, 8116 bytes, the same shared Stacks-wide preamble the published chow and etale-cohomology cards use); the first line of the pinned chapter file is \textbackslash{}input\{preamble\}, so no preamble text lives inside the chapter itself. The card takes that file's lines 8-18 (the verbatim package and the three \textbackslash{}comment-based environments reference/slogan/history at lines 9-12, section 13.2, together with the Xy-pic setup) and its lines 154 to the runtime end of the file (the native theorem declarations and mathematical macros). Not one character was changed, no package was added, no macro was defined, and the excerpt serves only as preamble, never as a span coordinate.\par\medskip
\noindent\textit{Editorial note.} The same-chapter pointers are located at their \textbackslash{}label lines inside the pinned chapter file; the cross-chapter pointers are located in the project's own tag list staged in the snapshot (sources/stacks-chow/tags.txt), at the verbatim lines carrying their canonical tags.\par\medskip
\noindent\textit{Editorial note.} The counted proof refers to the definition of the perfect closure and to same-chapter statements; the definition is recorded with kind definition in the navigation stubs.\par\medskip
\noindent\textit{Editorial note.} The counted spans contain no cite, so the card carries no bibliography.\par\medskip
\noindent\textit{Allowed standard prerequisites (editorial).}\begin{itemize}
\item The perfect closure of a separably closed field. Used at the start of the counted proof (lines 11330-11332). Turns the separably closed hypothesis into an algebraically closed coefficient field.
\item Irreducibility of fibres for a flat ring map. Invoked at the end of the counted proof (lines 11356-11358). Transfers irreducibility from the fibres to the total space.
\end{itemize}\medskip
\ifdefined\hypertarget\hypertarget{sourcepointer-1}{}\fi
\noindent\textit{Source reference.} Original definition 046X: sources/stacks-algebra/pinned-original-source.tex, source lines 10786--10786. Navigation pointer for \textbackslash{}ref\{definition-perfection\} used inside the counted span (line 11331); the argument that consumes it is retained in the span.\par\medskip
\ifdefined\hypertarget\hypertarget{sourcepointer-2}{}\fi
\noindent\textit{Source reference.} Original lemma 00G1: sources/stacks-algebra/pinned-original-source.tex, source lines 6885--6885. Navigation pointer for \textbackslash{}ref\{lemma-finite-type-field-Jacobson\} used inside the counted span (line 11356); the argument that consumes it is retained in the span.\par\medskip
\ifdefined\hypertarget\hypertarget{sourcepointer-3}{}\fi
\noindent\textit{Source reference.} Original lemma 00I6: sources/stacks-algebra/pinned-original-source.tex, source lines 11300--11300. Navigation pointer for \textbackslash{}ref\{lemma-flat-fibres-irreducible\} used inside the counted span (line 11357); the argument that consumes it is retained in the span.\par\medskip
\ifdefined\hypertarget\hypertarget{sourcepointer-4}{}\fi
\noindent\textit{Source reference.} Original lemma 00I3: sources/stacks-algebra/pinned-original-source.tex, source lines 10323--10323. Navigation pointer for \textbackslash{}ref\{lemma-limit-argument\} used inside the counted span (line 11345); the argument that consumes it is retained in the span.\par\medskip
\ifdefined\hypertarget\hypertarget{sourcepointer-5}{}\fi
\noindent\textit{Source reference.} Original lemma 0BRA: sources/stacks-algebra/pinned-original-source.tex, source lines 11085--11085. Navigation pointer for \textbackslash{}ref\{lemma-p-ring-map\} used inside the counted span (line 11337); the argument that consumes it is retained in the span.\par\medskip
\ifdefined\hypertarget\hypertarget{sourcepointer-6}{}\fi
\noindent\textit{Source reference.} Original lemma 00I4: sources/stacks-algebra/pinned-original-source.tex, source lines 10798--10798. Navigation pointer for \textbackslash{}ref\{lemma-perfect-reduced\} used inside the counted span (line 11343); the argument that consumes it is retained in the span.\par\medskip
\ifdefined\hypertarget\hypertarget{sourcepointer-7}{}\fi
\noindent\textit{Source reference.} Original lemma 00FV: sources/stacks-algebra/pinned-original-source.tex, source lines 6741--6741. Navigation pointer for \textbackslash{}ref\{theorem-nullstellensatz\} used inside the counted span (line 11354); the argument that consumes it is retained in the span.\par\medskip
\begin{lemma}
\label{lemma-separably-closed-irreducible}
Let $k$ be a separably closed field.
Let $R$, $S$ be $k$-algebras. If $R$, $S$ have a unique
minimal prime, so does $R \otimes_k S$.
\end{lemma}

\begin{proof}
Let $k \subset \overline{k}$ be a perfect closure, see
Definition \ref{definition-perfection}.
By assumption $\overline{k}$ is algebraically closed.
The ring maps $R \to R \otimes_k \overline{k}$ and
$S \to S \otimes_k \overline{k}$ and
$R \otimes_k S \to (R \otimes_k S) \otimes_k \overline{k}
= (R \otimes_k \overline{k}) \otimes_{\overline{k}} (S \otimes_k \overline{k})$
satisfy the assumptions of Lemma \ref{lemma-p-ring-map}.
Hence we may assume $k$ is algebraically closed.

\medskip\noindent
We may replace $R$ and $S$ by their reductions.
Hence we may assume that $R$ and $S$ are domains.
By Lemma \ref{lemma-perfect-reduced} we see that $R \otimes_k S$ is
reduced. Hence its spectrum is reducible if and only if it contains a nonzero
zerodivisor. By Lemma \ref{lemma-limit-argument} we reduce to the case where
$R$ and $S$ are domains of finite type over $k$ algebraically closed.

\medskip\noindent
Note that the ring map $R \to R \otimes_k S$ is of finite
presentation and flat. Moreover, for every maximal ideal
$\mathfrak m$ of $R$ we have
$(R \otimes_k S) \otimes_R R/\mathfrak m \cong S$ because
$k \cong R/\mathfrak m$ by the Hilbert Nullstellensatz Theorem
\ref{theorem-nullstellensatz}. Moreover, the set of
maximal ideals is dense in the spectrum of $R$ since
$\Spec(R)$ is Jacobson, see Lemma \ref{lemma-finite-type-field-Jacobson}.
Hence we see that Lemma \ref{lemma-flat-fibres-irreducible} applies
to the ring map $R \to R \otimes_k S$ and we conclude that
the spectrum of $R \otimes_k S$ is irreducible as desired.
\end{proof}

\end{document}
